\documentclass[UTF8,fontset=none,11pt,a4paper]{ctexart}
\usepackage{ipara-handbook}
\handbooksetup{DS-B01 Frontier Traversal}{408 · 数据结构 Bridge B01}{Visited · Frontier · Expansion}
\begin{document}
\handbooktitle{Frontier Traversal}{递归、栈、队列怎样统一成“待展开边界”}{408 · 数据结构 Internal Bridge B01}

\begin{corebox}{桥梁接口}
\[
\text{Visited State}+\text{Frontier}\to\text{Choose Next}\to\text{Visit}\to\text{Expand}\to\text{Update Frontier}.
\]
递归调用栈、显式栈和队列都保存已发现但尚未展开的对象；它们的差异在取出顺序，而不是把 Tree、Graph、Stack、Queue 的本体重新定义一遍。
\end{corebox}
\begin{methodbox}{Owners 与停止边界}
\begin{itemize}
\item DS02 Owns LIFO/FIFO 端点；DS04 Owns 树节点与遍历访问时机；DS07 Owns 图表示、visited 与覆盖。
\item 本桥只 Owns 可复用的 frontier 接口、发现时机和覆盖检查。
\item 树无环时不需要一般图的 visited；图不连通时必须由外层扫描补成遍历森林。
\end{itemize}
\end{methodbox}
\tableofcontents\newpage

\section{接口语义}
把 frontier 分为三种状态：未发现、已发现未展开、已展开。发现时加入 frontier 并立即标记，可保证每个对象最多进入边界一次；取出时才处理其邻接对象。若等到取出才标记，两个前驱可能重复入队，复杂度和输出都会改变。

\section{取出规则决定顺序}
栈是 LIFO，形成深度优先展开；队列是 FIFO，形成按层展开。递归 DFS 只是由语言调用栈隐式保存同一状态。对有权候选，优先队列也可作为 frontier，但此时取出的是当前最小优先级，正确性要转交 DS08 的目标不变量。
\begin{examplebox}{树与图的同一镜头}
树的先序/中序/后序改变“处理根”的时机；图的 DFS/BFS 改变 frontier 的取出顺序。树的无环结构允许省略 visited，图则必须显式维护它。
\end{examplebox}

\section{代码级最小调用}
\begin{lstlisting}[language=C++]
frontier.push(start); visited[start] = true;
while (!frontier.empty()) {
  auto u = frontier.take();
  visit(u);
  for (auto v : expand(u))
    if (!visited[v]) { visited[v] = true; frontier.push(v); }
}
\end{lstlisting}
具体栈/队列实现分别见 DS02，Tree/Graph 的展开函数见 DS04/DS07；本段只固定交接顺序。

\section{检查、复杂度与边界}
每个对象最多发现一次、每条邻接关系最多被检查一次，是邻接表遍历 $O(V+E)$ 的证明骨架。检查空 frontier、空树、越界起点、非连通分量和重复边；若题目要求最短权值或最小代价，不把 FIFO 的层数结论外推到加权图。

\section{做题控制协议}
先标出 frontier、visited 和 expand 的对象；再问取出规则；最后写覆盖终止条件。若序列不唯一，说明邻接顺序；若只覆盖一个分量，明确停止位置；若需要全图覆盖，补外层未访问顶点扫描。
\begin{corebox}{压缩复原}
忘记 DFS/BFS 口诀时，复原三问：边界保存什么？下一次取谁？发现是在入边界还是出边界时标记？
\end{corebox}
\end{document}
